\documentclass[11pt]{article}
\usepackage[letterpaper,margin=0.75in]{geometry}
\usepackage[T1]{fontenc}
\usepackage[utf8]{inputenc}
\usepackage{lmodern}
\usepackage{microtype}
\usepackage{xcolor}
\usepackage{enumitem}
\usepackage{booktabs}

\definecolor{titleline}{RGB}{131,44,39}
\setlength{\parindent}{0pt}
\setlength{\parskip}{0.55em}
\setlist[itemize]{leftmargin=1.2em,itemsep=0.25em,topsep=0.25em}
\setlist[enumerate]{leftmargin=1.4em,itemsep=0.35em,topsep=0.35em}

\begin{document}

{\LARGE\bfseries Scenario 3 Example of Play\par}
{\large Harper's Ferry to Port Republic, Turn 1\par}
\vspace{0.25em}
\color{titleline}\rule{\textwidth}{1.2pt}\par
\vspace{0.45em}

This example follows one plausible full turn in Scenario 3 and is meant to illustrate command, movement, reaction, combat, retreat, and conclusion procedures. The die rolls are chosen for teaching purposes. The map scan available is not crisp enough for a perfect hex-by-hex movement proof, so the example assumes Jackson advances along the main road corridor from hex 4511 toward Shields in hex 3519.

\textbf{Starting situation.}

\begin{itemize}
\item Confederate: Jackson, Winder, Ewell, Ashby, and a wagon begin in hex 4511.
\item Union: Banks, Williams, Brodhead, and Detachment E begin in hex 4502; Saxton is in 4611; Fremont, Blenker, and Bayard are in 2516; Shields, Kimball, and Detachment F are in 3519.
\item Scenario 3 is a 1862 scenario, so the Confederate player gains the initiative on a die roll of 1--5.
\end{itemize}

\textbf{1. Initial Phase}

There are no reinforcements and no replacements in this example turn, so the phase ends immediately.

\textbf{2. Initiative Phase}

The Confederate player rolls a \texttt{4}. In a 1862 scenario that gives the Confederates the initiative, so they will be the first player this turn.

\textbf{3. Confederate Command Phase}

Jackson is the only Confederate supreme commander on the map, so only one command-point die roll is needed. The Confederate player rolls a \texttt{5}. On the Command Point Table that gives Jackson \texttt{3} CPs for the turn.

Jackson's command is good enough that \texttt{3} CPs give him the full \texttt{100\%} of his printed movement allowance. Infantry and leaders normally have \texttt{8} MPs, so Jackson's force will be able to use the full infantry allowance before stacking effects.

\textbf{4. Confederate Movement Phase}

Jackson moves with Winder, Ewell, and Ashby as one force. Because more than two divisions are moving together, the force loses \texttt{1} MP for stacking under rule 7.31, so the moving force is treated as having \texttt{7} MPs.

The Confederate player moves the force out of hex 4511 along the road toward Shields in 3519.

\begin{enumerate}
\item Most of the move uses road movement, so the force advances efficiently down the valley.
\item Once Jackson enters a hex adjacent to 3519, the Union player checks reaction. Shields's stack is within reaction range, so the Union player declares a reaction attempt with Shields, Kimball, and Detachment F.
\item The Union player rolls for reaction and fails, so the stack remains in 3519.
\item Jackson then enters hex 3519. Because it contains enemy combat units, the Confederate force pays the extra \texttt{+1} MP required by rule 7.41 and may spend no more MPs this Movement Phase.
\end{enumerate}

At the end of Confederate movement, hex 3519 contains both sides, so a battle is now set up for the Combat Phase unless the Union can change the situation during its half of the turn.

\textbf{5. Union Command Phase}

The Union has three supreme commanders in the 1862 scenarios, so Banks, Fremont, and Shields each roll separately for command points.

\begin{itemize}
\item Banks rolls a \texttt{4} and receives \texttt{2} CPs.
\item Fremont rolls a \texttt{4} and also receives \texttt{2} CPs.
\item Shields rolls a \texttt{5} and receives \texttt{3} CPs.
\end{itemize}

Banks and Fremont therefore get only moderate movement this turn, while Shields's force in 3519 will have full movement available if it tries to leave the hex.

\textbf{6. Union Movement Phase}

The Union player now decides what to do about Jackson's incursion.

\begin{itemize}
\item Banks keeps his main force at 4502 to protect the depot area and nearby roads.
\item Saxton stays in 4611.
\item Fremont advances part of his force eastward, but the distance is too great for Blenker and Bayard to intervene in the 3519 fight this turn.
\item Shields could attempt to leave the enemy-occupied hex using rule 7.43, but the Union player decides not to. The stack is in entrenchments, has usable command points, and would rather fight from a doubled defensive position than give up 3519.
\end{itemize}

No other movement changes the impending battle, so the turn proceeds to combat.

\textbf{7. Combat Phase}

Only one battle must be resolved: hex 3519.

\textbf{Supply check.}

Both sides are supplied.

\begin{itemize}
\item The Confederate force can use the wagon stacked with Jackson.
\item The Union force can use the wagon in 3519.
\end{itemize}

\textbf{Combat options.}

The Confederate player chooses \textit{Pitched Battle}. The Union player chooses \textit{Stand}. On the Combat Options Chart, that means the battle will last \texttt{2} rounds.

\textbf{Wagon depletion.}

Because Jackson chose \textit{Pitched Battle}, his wagon may deplete on a die roll of \texttt{5} or \texttt{6}. The Confederate player rolls a \texttt{5}, so the wagon is marked depleted. The Union wagon does not deplete.

\textbf{Commitment.}

Each side now determines how much of its force can be committed.

\begin{itemize}
\item Jackson rolls \texttt{8} on 2d6, adds \texttt{+1} for \textit{Pitched Battle}, and commits \texttt{100\%} of his force.
\item Shields rolls \texttt{4} on 2d6. Because the Union force did not move and is entrenched, the roll is heavily improved, so \texttt{80\%} of the defending force is committed.
\end{itemize}

That gives:

\begin{itemize}
\item Confederate committed strength: Winder \texttt{24} + Ewell \texttt{22} + Ashby \texttt{3} = \texttt{49} SPs.
\item Union committed strength: Kimball \texttt{24} + Detachment F \texttt{2} = \texttt{26} SPs total, of which \texttt{80\%} means \texttt{21} SPs committed.
\item Because the Union is defending in entrenchments, its committed strength is doubled to \texttt{42} for fire purposes, and the Confederate fire column shifts one column left.
\end{itemize}

\textbf{Round 1.}

\begin{itemize}
\item Confederate fire: Jackson's \textit{Pitched Battle} and combat bonus produce a modified die roll of \texttt{7}. On the shifted column that inflicts \texttt{6} SP losses on the Union.
\item Union fire: the defender gets \texttt{+1} to fire and ends with a modified die roll of \texttt{5}. That inflicts \texttt{3} SP losses on the Confederates.
\end{itemize}

The losses are applied simultaneously. The Union has now lost more than \texttt{10\%} of its initial committed force, so it checks for demoralization at the end of the round.

\begin{itemize}
\item Union demoralization: the Union has lost \texttt{6} of \texttt{21} committed SPs, which is over \texttt{20\%}. The demoralization modifiers are therefore \texttt{+2} for losses and \texttt{-1} because the defender is entrenched. A sample die roll of \texttt{4} becomes \texttt{5}, so the Union holds.
\item Confederate demoralization: the Confederates have lost only \texttt{3} of \texttt{49} committed SPs, which is below the \texttt{10\%} threshold, so they do not check.
\end{itemize}

\textbf{Round 2.}

\begin{itemize}
\item Confederate fire: a modified die roll of \texttt{6} inflicts \texttt{5} more SP losses.
\item Union fire: a modified die roll of \texttt{4} inflicts \texttt{3} more SP losses.
\end{itemize}

After two rounds, the total losses are:

\begin{itemize}
\item Union: \texttt{11} SPs lost.
\item Confederate: \texttt{6} SPs lost.
\end{itemize}

Both sides now check demoralization.

\begin{itemize}
\item Union demoralization: \texttt{11} losses out of \texttt{21} committed SPs is more than \texttt{50\%}. Using the same modifiers as before, a die roll of \texttt{5} becomes a failure and the Union force is demoralized.
\item Confederate demoralization: \texttt{6} losses out of \texttt{49} committed SPs is just over \texttt{10\%}. A die roll of \texttt{2}, modified by losses and Jackson's strong command, is low enough that the Confederates remain in good order.
\end{itemize}

Because the Union becomes demoralized, the battle ends with a Confederate victory.

\textbf{8. Retreat}

The Union force in 3519 must retreat immediately. Under rule 14.21 it retreats \texttt{4} MPs without paying ZOC costs, and the Union player chooses the legal path that pulls the force back toward friendly support while avoiding prohibited enemy-occupied hexes and river restrictions.

Since the retreating side lost fewer than \texttt{15} SPs, the Confederate result is a \textbf{minor victory}, not a major victory.

Jackson remains in 3519 at the end of combat. His wagon remains with the force, but it is now depleted.

\textbf{9. Conclusion Phase}

\begin{itemize}
\item \textbf{Depletion and replenishment:} the Confederate wagon used in the pitched battle is flipped to its depleted side at the end of the Combat Phase.
\item \textbf{Morale Recovery:} the Union may attempt to recover the demoralized force in the Morale Recovery Phase. In this example the recovery roll fails, so the force remains demoralized into the next turn.
\item \textbf{Union Panic Track:} because this is a 1862 scenario, the players now apply the appropriate Panic Track adjustment for the turn's result.
\end{itemize}

\textbf{What this example showed.}

\begin{itemize}
\item Initiative in a 1862 scenario
\item Command-point rolls and movement percentages
\item Stacking penalties during movement
\item Reaction movement failure
\item Entering an enemy-occupied hex
\item Supply and wagon depletion
\item Combat options and commitment
\item Multi-round combat
\item Demoralization and retreat
\item End-of-turn morale recovery and depleted wagon status
\end{itemize}

\end{document}
